% GENERATED FILE. Edit canonical lessons, then run npm run generate:books.
% canonical-source-hash: fnv1a64:6d48a32d2c51e5c9
% canonical-lessons: PA-C296-khisakna, PA-C296-chipakna, PA-C296-garha, PA-C296-chikna, PA-C296-pighalna

\chapter{Slippery Paths}
\label{ch:pa-a2-slippery-paths}
\hlchaptermodality{296}

\begin{chapteropening}
\textbf{By the end of this chapter:} \emph{I can describe slippery and sticky things.}

Complete the last lesson of chapter 296: I can describe slippery and sticky things.
\end{chapteropening}

\section[\texorpdfstring{\pa{ਖਿਸਕਣਾ} (khisakṇā)}{ਖਿਸਕਣਾ (khisakṇā)}]{\pa{ਖਿਸਕਣਾ} (khisakṇā) — to slide; to slip away}

\label{lesson:PA-C296-khisakna}



\begin{quote}
Before the new one: say the Punjabi for a ravine, then the Punjabi for to rot.
\end{quote}

\subsection*{You'll want to know: \pa{ਖਿਸਕਣਾ}}
\textbf{\pa{ਖਿਸਕਣਾ}} — \emph{khisakṇā} — \textquotedblleft{}to slide; to slip away\textquotedblright{}.

\begin{quote}
\pa{ਮੀਂਹ} \pa{ਵਿੱਚ} \pa{ਪਹਾੜ} \pa{ਤੋਂ} \pa{ਪੱਥਰ} \pa{ਖਿਸਕਦੇ} \pa{ਹਨ।}
\end{quote}

\begin{quote}
\emph{mīṁh vich pahāṛ ton patthar khisakde han.}
\end{quote}

\begin{quote}
In the rain, stones slide down from the mountain.
\end{quote}

Past (it …): \pa{ਉਹ} \pa{ਪੱਥਰ} \pa{ਖਿਸਕਿਆ}, \emph{oh patthar khiskiā}.

\begin{cousinweb}
\textbf{\pa{ਖਿਸਕਣਾ}} (\emph{khisakṇā}) is inherited from a reconstructed \emph{*khis}, \textquotedblleft{}slip\textquotedblright{}; CDIAL lists Punjabi \emph{khiskaṇā}, \textquotedblleft{}to slip away\textquotedblright{}, under it (CDIAL 3888).
\end{cousinweb}

\begin{grammarlens}[title={What you've built}]
One more word for this chapter.
\end{grammarlens}

\subsection*{Your turn}
\begin{itemize}
\raggedright
  \item \emph{Say it:} \emph{khisakṇā}
  \item \emph{Say it:} \emph{khisakṇā}, once more
  \item \emph{Say it:} \emph{mīṁh vich pahāṛ ton patthar khisakde han.} — In the rain, stones slide down from the mountain.
  \item \emph{Recall:} say the Punjabi for a ravine, then the Punjabi for to rot, then say \emph{khisakṇā} again
\end{itemize}

\begin{tcolorbox}[breakable,colback=teal!4,colframe=teal!35!black,title={Before you move on}]
What does \pa{ਖਿਸਕਣਾ} mean? (\textquotedblleft{}To slide; to slip away\textquotedblright{}.) Say it once more.
\end{tcolorbox}
\section[\texorpdfstring{\pa{ਚਿਪਕਣਾ} (chipakṇā)}{ਚਿਪਕਣਾ (chipakṇā)}]{\pa{ਚਿਪਕਣਾ} (chipakṇā) — to stick; to cling}

\label{lesson:PA-C296-chipakna}



\begin{quote}
Before the new one: say the Punjabi for to rot, then the Punjabi for to slide.
\end{quote}

\subsection*{You'll want to know: \pa{ਚਿਪਕਣਾ}}
\textbf{\pa{ਚਿਪਕਣਾ}} — \emph{chipakṇā} — \textquotedblleft{}to stick; to cling\textquotedblright{}.

\begin{quote}
\pa{ਚਿੱਕੜ} \pa{ਜੁੱਤੀ} \pa{ਨਾਲ} \pa{ਚਿਪਕਦਾ} \pa{ਹੈ।}
\end{quote}

\begin{quote}
\emph{chikkaṛ juttī nāl chipakdā hai.}
\end{quote}

\begin{quote}
Mud sticks to the shoe.
\end{quote}

Past (it …): \pa{ਉਹ} \pa{ਪੱਤਾ} \pa{ਚਿਪਕਿਆ}, \emph{oh pattā chipkiā}.

\begin{cousinweb}
\textbf{\pa{ਚਿਪਕਣਾ}} (\emph{chipakṇā}) is inherited from a reconstructed \emph{*cippa}, \textquotedblleft{}gummy matter\textquotedblright{}; CDIAL lists Punjabi \emph{cipakṇā}, \textquotedblleft{}to stick\textquotedblright{}, under it (CDIAL 4819).
\end{cousinweb}

\begin{grammarlens}[title={What you've built}]
One more word for this chapter.
\end{grammarlens}

\subsection*{Your turn}
\begin{itemize}
\raggedright
  \item \emph{Say it:} \emph{chipakṇā}
  \item \emph{Say it:} \emph{chipakṇā}, once more
  \item \emph{Say it:} \emph{chikkaṛ juttī nāl chipakdā hai.} — Mud sticks to the shoe.
  \item \emph{Recall:} say the Punjabi for to rot, then the Punjabi for to slide, then say \emph{chipakṇā} again
  \item \emph{Write it:} copy \textbf{\pa{ਚਿਪਕਣਾ}} once, checking the whole word
\end{itemize}

\begin{tcolorbox}[breakable,colback=teal!4,colframe=teal!35!black,title={Before you move on}]
What does \pa{ਚਿਪਕਣਾ} mean? (\textquotedblleft{}To stick; to cling\textquotedblright{}.) Say it once more.
\end{tcolorbox}
\section[\texorpdfstring{\pa{ਗਾੜ੍ਹਾ} (gāṛhā)}{ਗਾੜ੍ਹਾ (gāṛhā)}]{\pa{ਗਾੜ੍ਹਾ} (gāṛhā) — thick (of a liquid)}

\label{lesson:PA-C296-garha}



\begin{quote}
Before the new one: say the Punjabi for to slide, then the Punjabi for to stick.
\end{quote}

\subsection*{You'll want to know: \pa{ਗਾੜ੍ਹਾ}}
\textbf{\pa{ਗਾੜ੍ਹਾ}} — \emph{gāṛhā} — \textquotedblleft{}thick (of a liquid)\textquotedblright{}.

\begin{quote}
\pa{ਇਹ} \pa{ਦੁੱਧ} \pa{ਬਹੁਤ} \pa{ਗਾੜ੍ਹਾ} \pa{ਹੈ।}
\end{quote}

\begin{quote}
\emph{eh dudh bahut gāṛhā hai.}
\end{quote}

\begin{quote}
This milk is very thick.
\end{quote}

\begin{cousinweb}
\textbf{\pa{ਗਾੜ੍ਹਾ}} (\emph{gāṛhā}) is inherited from Sanskrit \emph{gāḍha}, \textquotedblleft{}pressed together, thick\textquotedblright{}; CDIAL lists Punjabi \emph{gāṛhā}, \textquotedblleft{}thick, close\textquotedblright{}, under it (CDIAL 4118).
\end{cousinweb}

\begin{grammarlens}[title={What you've built}]
One more word for this chapter.
\end{grammarlens}

\subsection*{Your turn}
\begin{itemize}
\raggedright
  \item \emph{Say it:} \emph{gāṛhā}
  \item \emph{Say it:} \emph{gāṛhā}, once more
  \item \emph{Say it:} \emph{eh dudh bahut gāṛhā hai.} — This milk is very thick.
  \item \emph{Recall:} say the Punjabi for to slide, then the Punjabi for to stick, then say \emph{gāṛhā} again
  \item \emph{Write it:} copy \textbf{\pa{ਗਾੜ੍ਹਾ}} once, checking the whole word
\end{itemize}

\begin{tcolorbox}[breakable,colback=teal!4,colframe=teal!35!black,title={Before you move on}]
What does \pa{ਗਾੜ੍ਹਾ} mean? (\textquotedblleft{}Thick (of a liquid)\textquotedblright{}.) Say it once more.
\end{tcolorbox}
\section[\texorpdfstring{\pa{ਚਿਕਨਾ} (chiknā)}{ਚਿਕਨਾ (chiknā)}]{\pa{ਚਿਕਨਾ} (chiknā) — slick; greasy}

\label{lesson:PA-C296-chikna}



\begin{quote}
Before the new one: say the Punjabi for to stick, then the Punjabi for thick.
\end{quote}

\subsection*{You'll want to know: \pa{ਚਿਕਨਾ}}
\textbf{\pa{ਚਿਕਨਾ}} — \emph{chiknā} — \textquotedblleft{}slick; greasy\textquotedblright{}.

\begin{quote}
\pa{ਮੀਂਹ} \pa{ਤੋਂ} \pa{ਬਾਅਦ} \pa{ਫ਼ਰਸ਼} \pa{ਚਿਕਨਾ} \pa{ਹੈ।}
\end{quote}

\begin{quote}
\emph{mīṁh ton bāad farsh chiknā hai.}
\end{quote}

\begin{quote}
After the rain the floor is slick.
\end{quote}

\begin{cousinweb}
\textbf{\pa{ਚਿਕਨਾ}} (\emph{chiknā}) is inherited from Sanskrit \emph{cikkaṇa}, \textquotedblleft{}slippery, unctuous\textquotedblright{}; CDIAL lists Punjabi \emph{ciknā}, \textquotedblleft{}sticky, oily, smooth\textquotedblright{}, under it (CDIAL 4782).
\end{cousinweb}

\begin{grammarlens}[title={What you've built}]
One more word for this chapter.
\end{grammarlens}

\subsection*{Your turn}
\begin{itemize}
\raggedright
  \item \emph{Say it:} \emph{chiknā}
  \item \emph{Say it:} \emph{chiknā}, once more
  \item \emph{Say it:} \emph{mīṁh ton bāad farsh chiknā hai.} — After the rain the floor is slick.
  \item \emph{Recall:} say the Punjabi for to stick, then the Punjabi for thick, then say \emph{chiknā} again
  \item \emph{Write it:} copy \textbf{\pa{ਚਿਕਨਾ}} once, checking the whole word
\end{itemize}

\begin{tcolorbox}[breakable,colback=teal!4,colframe=teal!35!black,title={Before you move on}]
What does \pa{ਚਿਕਨਾ} mean? (\textquotedblleft{}Slick; greasy\textquotedblright{}.) Say it once more.
\end{tcolorbox}
\section[\texorpdfstring{\pa{ਪਿਘਲਣਾ} (pighalṇā)}{ਪਿਘਲਣਾ (pighalṇā)}]{\pa{ਪਿਘਲਣਾ} (pighalṇā) — to melt}

\label{lesson:PA-C296-pighalna}



\begin{quote}
Before the new one: say the Punjabi for thick, then the Punjabi for slick.
\end{quote}

\subsection*{You'll want to know: \pa{ਪਿਘਲਣਾ}}
\textbf{\pa{ਪਿਘਲਣਾ}} — \emph{pighalṇā} — \textquotedblleft{}to melt\textquotedblright{}.

\begin{quote}
\pa{ਧੁੱਪ} \pa{ਵਿੱਚ} \pa{ਬਰਫ਼} \pa{ਪਿਘਲਦੀ} \pa{ਹੈ।}
\end{quote}

\begin{quote}
\emph{dhupp vich baraf pighaldī hai.}
\end{quote}

\begin{quote}
Ice melts in the sunshine.
\end{quote}

Past (it …): \pa{ਉਹ} \pa{ਬਰਫ਼} \pa{ਪਿਘਲੀ}, \emph{oh baraf pighlī}.

\begin{cousinweb}
\textbf{\pa{ਪਿਘਲਣਾ}} (\emph{pighalṇā}) is inherited from a reconstructed \emph{*praghirati}, \textquotedblleft{}flows forth\textquotedblright{}; CDIAL lists Punjabi \emph{pighalnā}, \textquotedblleft{}to be dissolved, be melted\textquotedblright{}, under it (CDIAL 8486).
\end{cousinweb}

\begin{grammarlens}[title={What you've built}]
One more word for this chapter.
\end{grammarlens}

\subsection*{Your turn}
\begin{itemize}
\raggedright
  \item \emph{Say it:} \emph{pighalṇā}
  \item \emph{Say it:} \emph{pighalṇā}, once more
  \item \emph{Say it:} \emph{dhupp vich baraf pighaldī hai.} — Ice melts in the sunshine.
  \item \emph{Recall:} say the Punjabi for thick, then the Punjabi for slick, then say \emph{pighalṇā} again
  \item \emph{Write it:} copy \textbf{\pa{ਪਿਘਲਣਾ}} once, checking the whole word
\end{itemize}

\begin{tcolorbox}[breakable,colback=teal!4,colframe=teal!35!black,title={Before you move on}]
What does \pa{ਪਿਘਲਣਾ} mean? (\textquotedblleft{}To melt\textquotedblright{}.) Say it once more.
\end{tcolorbox}
